\chapter{The Physical Layer and Cable Faults}\label{ch:physical}
\bp{1.1}

\begin{outcomes}
By the end of this chapter you will be able to:
\begin{tight}
  \item \textbf{Select} the right medium --- copper, multimode fibre, single-mode
        fibre --- for a given distance and speed.
  \item \textbf{Read} the interface counters that expose a physical fault, and
        say which counter implicates which cause.
  \item \textbf{Diagnose} the classic faults: duplex mismatch, speed mismatch,
        wrong cable type, bad pin-out, dirty or over-long fibre, and a
        transceiver operating outside its optical budget.
  \item \textbf{Distinguish} a port that is down from a port that is up and
        damaging traffic --- the second is far more dangerous.
\end{tight}
\end{outcomes}

\begin{prereq}
\chref{models} for the layered model and the idea of a collision domain.
\end{prereq}

\begin{keyterms}
twisted pair $\bullet$ Cat5e/Cat6 $\bullet$ multimode $\bullet$ single-mode
$\bullet$ SFP $\bullet$ attenuation $\bullet$ auto-negotiation $\bullet$ duplex
mismatch $\bullet$ CRC error $\bullet$ runt $\bullet$ giant $\bullet$ late
collision $\bullet$ input error $\bullet$ pin-out $\bullet$ straight-through
$\bullet$ crossover $\bullet$ Auto-MDIX
\end{keyterms}

\begin{examnote}[Read the topic wording]
Topic \tp{1.1} says \emph{diagnose} interface and cable issues, and then names
them: ``collisions, errors, mismatched duplex, speed, distance, interface, signal
levels, pin out, and cable types''. That list is a syllabus in itself. You will
not be asked to describe a cable; you will be shown counters and asked what is
wrong.
\end{examnote}

\section{Choosing the medium}

\begin{definitionbox}[Bits per second, and why it is not bytes]
Link speeds are always quoted in \textbf{bits} per second: 1~Gbps is one thousand
million \emph{bits} each second. File sizes are quoted in \textbf{bytes}, and a
byte is eight bits.

So a 1~Gbps link moves at most about 125~megabytes per second, not 1000 --- divide
by eight. The lower-case \cmd{b} means bits and the upper-case \cmd{B} means
bytes, which is the entire difference between \cmd{Mbps} and \cmd{MBps}.

Two further distinctions worth having straight from the start:
\begin{tight}
  \item \term{Bandwidth} is the rated capacity of the link --- what it could carry.
  \item \term{Throughput} is what you actually got, after overhead, contention,
        errors and whatever the far end could manage.
\end{tight}
Throughput is always lower, and a user reporting ``slow'' is reporting
throughput. Nothing in this chapter changes bandwidth; several things in it
quietly destroy throughput.
\end{definitionbox}

\begin{definitionbox}[Duplex, and attenuation]
\term{Duplex} is who may transmit, and when.

\begin{tight}
  \item \textbf{Half duplex} --- one end at a time, like a walkie-talkie. If both
        transmit at once the signals overlap and the frames are destroyed; that
        overlap is a \term{collision}, and each end is expected to detect it,
        back off and retry.
  \item \textbf{Full duplex} --- both ends at once, on separate pairs of wires,
        like a telephone call. \textbf{Collisions are impossible}, so a
        full-duplex port that reports collisions is telling you something is
        badly wrong. That is the single most useful sentence in this chapter.
\end{tight}

Modern switched links are full duplex, and both ends normally agree on it through
\term{auto-negotiation} --- a brief exchange when the link comes up in which each
end states what it can do and both pick the best they share.

\term{Attenuation} is separate and simpler: a signal loses strength with
distance and through material. Every distance limit in the table below is really
an attenuation limit --- the point at which the far end can no longer reliably
tell a 1 from a 0.
\end{definitionbox}

\begin{longtable}{@{}L{3.1cm}L{2.4cm}L{2.6cm}L{6.8cm}@{}}
\toprule
\rowcolor{primary!15}
\textbf{Medium} & \textbf{Typical reach} & \textbf{Speed} & \textbf{Use it when} \\
\midrule
\endhead
Cat5e twisted pair & 100~m & 1~Gbps & Desk to access switch. The default, and
the distance limit is a hard 100~m including patch leads \\
\rowcolor{lightbg}
Cat6 / Cat6a & 100~m (55~m for 10G on Cat6) & 1--10~Gbps & Uplinks within a
building, or any run you expect to upgrade \\
Multimode fibre (OM3/OM4) & 300--400~m & 1--10~Gbps & Between floors or between
buildings on one campus. Cheaper optics, shorter reach \\
\rowcolor{lightbg}
Single-mode fibre & Kilometres & 1--100~Gbps & Between sites. More expensive
optics, effectively unlimited campus reach \\
Direct-attach copper & 1--10~m & 10~Gbps+ & Switch to switch inside one rack \\
\bottomrule
\end{longtable}

\begin{conceptbox}[Distance is a fault, not a preference]
The exam names \emph{distance} as a diagnosable issue for a reason. Exceed the
reach of a medium and the link does not fail cleanly --- it comes up and then
corrupts frames, which shows as climbing CRC errors rather than a down port. A
95~m run that works in winter and fails in summer is a real phenomenon, and the
answer is fibre, not another patch lead.
\end{conceptbox}

\subsection{Copper pin-outs, and why you can mostly forget them}

A \term{straight-through} cable connects pin 1 to pin 1; a \term{crossover}
swaps the transmit and receive pairs. Historically you needed a crossover between
two devices of the same kind --- switch to switch, PC to router --- and a
straight-through between unlike devices.

\term{Auto-MDIX} makes the switch detect and correct this itself, and it is on by
default on every modern platform. It is still examinable because it fails in two
situations: when a port has been hard-coded to a speed and duplex (Auto-MDIX
requires auto-negotiation to be on), and on older or simulated gear. If a link
between two switches will not come up and both ends are hard-coded, the pin-out
is a live suspect.

\section{The counters that name the fault}

\begin{verify}{SW1}{show interfaces GigabitEthernet0/1}
GigabitEthernet0/1 is up, line protocol is up (connected)
  Hardware is Gigabit Ethernet, address is 00d0.58a1.0201
  MTU 1500 bytes, BW 1000000 Kbit, DLY 10 usec
  Full-duplex, 1000Mb/s, media type is 10/100/1000BaseTX
  5 minute input rate 41000 bits/sec, 12 packets/sec
     284517 packets input, 39281044 bytes, 0 no buffer
     Received 1204 broadcasts (0 multicasts)
     0 runts, 0 giants, 0 throttles
     0 input errors, 0 CRC, 0 frame, 0 overrun, 0 ignored
     271033 packets output, 33771201 bytes, 0 underruns
     0 output errors, 0 collisions, 0 interface resets
     0 late collision, 0 deferred
\end{verify}

That is a healthy port: every error counter zero, full duplex, and the speed
matching the media. Learn what this looks like, because you diagnose by
difference.

\begin{longtable}{@{}L{3.4cm}L{5.4cm}L{6.1cm}@{}}
\toprule
\rowcolor{primary!15}
\textbf{Counter climbing} & \textbf{What it means} & \textbf{Most likely cause} \\
\midrule
\endhead
\cmd{CRC} and \cmd{input errors} & Frames arrived damaged & Bad cable, exceeded
distance, dirty fibre connector, electrical interference. If \cmd{collisions} is
zero, suspect the cable before the duplex \\
\rowcolor{lightbg}
\cmd{collisions} on a full-duplex link & Should be impossible & \textbf{Duplex
mismatch.} One end is half, the other full \\
\cmd{late collision} & A collision after the first 64 bytes & Duplex mismatch, or
a segment longer than the standard allows \\
\rowcolor{lightbg}
\cmd{runts} & Frames shorter than 64 bytes & Collisions, or a failing NIC \\
\cmd{giants} & Frames longer than the MTU & MTU mismatch, or a misconfigured
jumbo-frame setting on one side \\
\rowcolor{lightbg}
\cmd{input errors} with \cmd{CRC} zero & Damaged, but not by corruption & Often
an MTU or framing problem rather than the cable \\
\cmd{output drops} & The device could not transmit fast enough & Congestion, not
a cable fault. Do not replace the cable \\
\bottomrule
\end{longtable}

\begin{conceptbox}[The duplex mismatch, in one paragraph]
Auto-negotiation lets two ports agree on speed and duplex. If one end is set to
\cmd{duplex full} by hand and the other is left on auto, the auto end cannot
detect duplex, falls back to half, and you have a mismatch. The half-duplex end
listens before transmitting and treats simultaneous transmission as a collision;
the full-duplex end transmits whenever it likes. The link stays \emph{up}. It
just loses a percentage of every conversation, and the symptom the user reports
is ``the network is slow'', not ``the network is down''.

The fix is to set both ends the same way. The rule of thumb is: leave both on
auto, or hard-code both. Never one of each.
\end{conceptbox}

\begin{config}{a hard-coded access port, both ends matching}
interface GigabitEthernet0/1
 description Link to SW2 - speed and duplex fixed at both ends
 speed 1000
 duplex full
 no shutdown
\end{config}

\subsection{Fibre and signal levels}

Topic \tp{1.1} names \emph{signal levels}. On a fibre link the transceiver
reports how much light it is transmitting and receiving, in dBm, and a receive
level below the optic's sensitivity produces exactly the same symptom as a bad
copper cable --- a link that is up with climbing CRC errors.

\begin{verify}{SW1}{show interfaces GigabitEthernet0/25 transceiver}
                             Optical   Optical
           Temperature  Voltage  Tx Power  Rx Power
Port       (Celsius)    (Volts)  (dBm)     (dBm)
---------  -----------  -------  --------  --------
Gi0/25       33.2         3.31     -2.4     -28.9
\end{verify}

A receive power of $-28.9$~dBm against a typical sensitivity of about $-23$~dBm
means the link is starved of light. Causes, in the order worth checking: a dirty
or unseated connector, a bend or break in the fibre, too many patch panels, a
run longer than the optic supports, or a single-mode optic on multimode fibre.

\begin{pitfall}
\begin{tight}
  \item \textbf{Replacing the cable first.} It is the most expensive check in
        time and the least likely single cause on a link that has been working.
        Read the counters first: they tell you which of six causes it is.
  \item \textbf{Treating ``up/up'' as healthy.} A port can be up, passing
        traffic, and corrupting five per cent of it. The exam will show you
        exactly this and expect you not to declare it fine.
  \item \textbf{Hard-coding one end.} The commonest way to create a duplex
        mismatch is to fix one side ``to be safe''.
  \item \textbf{Reading raw totals instead of a rate.} Counters accumulate since
        the last clear. Ten thousand CRC errors over two years is noise; ten
        thousand since yesterday is a fault. Use \cmd{clear counters} and watch.
\end{tight}
\end{pitfall}

\begin{breakfix}[``The file server is slow, but it is not down''.]
A user copies a large file and gets a fraction of the expected throughput. Ping
succeeds with no loss. You look at the access port:

\begin{verify}{SW1}{show interfaces GigabitEthernet0/7 | include duplex|collision|CRC}
  Half-duplex, 100Mb/s, media type is 10/100/1000BaseTX
     14432 input errors, 118 CRC, 0 frame, 0 overrun, 0 ignored
     0 output errors, 89211 collisions, 0 interface resets
     2244 late collision, 0 deferred
\end{verify}

\textbf{Diagnosis.} The port is half duplex at 100~Mbps on gigabit-capable media,
with a very large collision count and --- decisively --- \emph{late} collisions.
Late collisions on a modern switched link are the signature of a duplex mismatch:
the far end is transmitting full duplex while this end still listens before
sending. The small CRC count is a consequence of the collisions, not a separate
cable fault.

\textbf{Fix.} Find the server NIC setting. It has almost certainly been hard-coded
to full duplex while the switch port was left on auto. Set both ends to auto, or
both to \cmd{speed 1000} / \cmd{duplex full}, then \cmd{clear counters} and
confirm the errors stop climbing. If they do not, \emph{now} suspect the cable.
\end{breakfix}

\begin{lab}{Reading a broken link}{Packet Tracer}
\textbf{Topology.} Two 2960 switches connected by one link; a PC on each.

\begin{tasks}
  \item Cable \cmd{SW1 Gi0/1} to \cmd{SW2 Gi0/1}. Confirm both PCs can ping each
        other, then run \cmd{show interfaces Gi0/1} on both and record the speed,
        duplex and every error counter while the link is healthy.
  \item On \cmd{SW1} only, set \cmd{speed 100} and \cmd{duplex half}. Leave
        \cmd{SW2} on auto.
  \item Generate traffic between the PCs. Record what changes on each side.
        Which counter moved first?
  \item Repeat with \cmd{SW1} at \cmd{duplex full} and \cmd{SW2} on auto ---
        the classic mismatch. Note that the link stays up.
  \item \textbf{Extension.} Without looking at your own configuration, hand the
        topology to a partner, have them introduce one fault from this chapter,
        and diagnose it from \cmd{show} output alone in under three minutes.
\end{tasks}

\textbf{Verification.} You should be able to state, from counters only, whether
a link is healthy, mismatched, or physically damaged.
\end{lab}

\begin{checkpoint}
\begin{tightnum}
  \item A gigabit link is up. \cmd{collisions} is increasing and the interface
        reports full duplex. What single conclusion can you draw, and why is it
        certain?
  \item A fibre link shows Rx power of $-31$~dBm against a sensitivity of
        $-23$~dBm. Give three causes in the order you would check them.
  \item Why is a port that is administratively down easier to deal with than one
        that is up with rising CRC errors?
\end{tightnum}
\end{checkpoint}

\begin{chaptersummary}
\begin{tight}
  \item Copper is 100~m, full stop. Multimode reaches a campus; single-mode
        reaches a city. Exceeding a medium's reach corrupts frames rather than
        dropping the link.
  \item Diagnose from counters, not from intuition. CRC without collisions points
        at the cable; collisions on a full-duplex link point at duplex; late
        collisions are close to proof of a duplex mismatch.
  \item Set both ends of a link the same way --- both auto, or both fixed.
  \item ``Up/up'' does not mean healthy. The dangerous fault is the one that
        passes most traffic and damages the rest.
  \item On fibre, read \cmd{show interfaces \ldots{} transceiver}. Light levels
        are a diagnosable fault the blueprint names explicitly.
\end{tight}
\end{chaptersummary}

\begin{reviewq}
  \item \textbf{(MCQ)} Which counter most strongly indicates a duplex mismatch?
        (a)~giants (b)~late collision (c)~output drops (d)~no buffer
  \item \textbf{(MCQ)} A 10~Gbps link is required over 2~km between two
        buildings. Which medium?
        (a)~Cat6a (b)~multimode OM4 (c)~single-mode fibre (d)~direct-attach copper
  \item \textbf{(Short)} Explain why hard-coding duplex on one end of a link and
        leaving the other on auto produces a mismatch rather than a failure to
        link.
  \item \textbf{(Short)} A port shows 8,000 CRC errors and zero collisions.
        State the two most likely causes and the command you would run next.
  \item \textbf{(Applied)} You inherit a network where every access port has been
        hard-coded to \cmd{speed 100} and \cmd{duplex full}. Describe the risk
        this creates, and the safest order in which to migrate the estate back to
        auto-negotiation.
\end{reviewq}
